\subsection{Stokes' formula for locally polyhedral sets}

In this number, X denotes a pure real manifold of finite dimension n of class \(C^{r}\) ( \(r \geqslant 2\) ); we assume X is separated.\footnote{The reader interested in more general cases may consult H. Whitney, \emph{Geometric Integration Theory}, Chap. III, § 18 (Princeton Univ. Press, 1957).}

11.3.1. Let A be a subset of a finite-dimensional real vector space. We say that A is polyhedral if it is a finite union of finite intersections of closed half-spaces.

A subset A of X is said to be locally polyhedral if, for every \(x \in X\), there exists a chart \(c = (\mathrm{U}, \varphi, \mathrm{E})\) of X at x and a polyhedral subset \(A_{c}\) of E such that \(\varphi(\mathrm{A} \cap \mathrm{U}) = \varphi(\mathrm{U}) \cap \mathrm{A}_{c}\). A piece of X is a locally polyhedral subset.

11.3.2. Let A be a closed subset of X, and let \(\mathrm{Fr}(\mathrm{A}) = \mathrm{A} - \mathring{\mathrm{A}}\) be its frontier. A point \(x \in \mathrm{Fr}(\mathrm{A})\) is said to be regular if there exists an open neighborhood U of x such that \(A \cap U\) is a piece of the manifold U (in which case x belongs to the boundary of \(A \cap U\)). We denote by \(\partial A\) the set of regular points of \(\mathrm{Fr}(\mathrm{A})\) and call it the regular boundary (or simply the boundary) of A. The set \(A' = \mathring{A} \cup \partial A\) is a piece of X, with boundary \(\partial A\).

11.3.3. * Example. --- Let \(\mathcal{H}\) be a locally finite set of hyperplanes of a finite-dimensional real affine space E of dimension \(n\), and let C be a chamber of E relative to \(\mathcal{H}\) (LIE, V, § 1, no. 3). The closure \(\overline{\mathbf{C}}\) of C is a locally polyhedral subset of the manifold E; the regular boundary of \(\overline{\mathbf{C}}\) is the union of the faces of C (loc. cit., no. 4). In particular, if C is an open simplex (loc. cit., no. 6) with vertices \(a_0, \ldots, a_n\), the boundary of \(\overline{\mathbf{C}}\) is the union of the \(\overline{\mathbf{C}_{\{i\}}}\) (\(0 \leqslant i \leqslant n\)), where \(\mathbf{C}_{\{i\}}\) is the open simplex with vertices \(a_0, \ldots, a_{i-1}, a_{i+1}, \ldots, a_n\) in the affine space generated by these vertices.* 
11.3.4. Let A be a locally polyhedral subset of X, and let \(\omega\) be a twisted differential form of degree \(n - 1\) on X with values in a Banach space E; suppose that \(\omega\) is of class \(\mathbf{C}^1\) and that the intersection of its support with A is compact. Then, the characteristic function of A (resp. \(A'\), \(\mathring{A}\)) is essentially integrable for the vector measure defined by \(d\omega\) on X, and one has

\[
\int_ {\mathrm{A}} d \omega = \int_ {\mathrm{A} ^ {\prime}} d \omega = \int_ {\mathring{\mathrm{A}}} d \omega.
\]

The differential form \(\omega|\partial\mathbf{A}\) (for the canonical orientation of the canonical injection of \(\partial\mathbf{A}\), considered as the boundary of the piece \(\mathbf{A}'\), into \(\mathbf{X}\)) is integrable on \(\partial\mathbf{A}\), and one has

\[
\int_ {\mathbf {A}} d \omega = \int_ {\partial \mathbf {A}} \omega \quad (\text{Stokes' formula}).
\]

When X is equipped with an orientation, and \(\partial\mathbf{A}\) is equipped with the corresponding orientation (11.2.1), this formula remains valid if \(\omega\) is an ordinary differential form of degree \(n-1\) with values in E, of class \(\mathbf{C}^{1}\), and such that \(\mathbf{A}\cap\operatorname{Supp}\omega\) is compact.

